% Options for packages loaded elsewhere
\PassOptionsToPackage{unicode}{hyperref}
\PassOptionsToPackage{hyphens}{url}
%
\documentclass[
]{article}
\usepackage{amsmath,amssymb}
\usepackage{iftex}
\ifPDFTeX
  \usepackage[T1]{fontenc}
  \usepackage[utf8]{inputenc}
  \usepackage{textcomp} % provide euro and other symbols
\else % if luatex or xetex
  \usepackage{unicode-math} % this also loads fontspec
  \defaultfontfeatures{Scale=MatchLowercase}
  \defaultfontfeatures[\rmfamily]{Ligatures=TeX,Scale=1}
\fi
\usepackage{lmodern}
\ifPDFTeX\else
  % xetex/luatex font selection
\fi
% Use upquote if available, for straight quotes in verbatim environments
\IfFileExists{upquote.sty}{\usepackage{upquote}}{}
\IfFileExists{microtype.sty}{% use microtype if available
  \usepackage[]{microtype}
  \UseMicrotypeSet[protrusion]{basicmath} % disable protrusion for tt fonts
}{}
\makeatletter
\@ifundefined{KOMAClassName}{% if non-KOMA class
  \IfFileExists{parskip.sty}{%
    \usepackage{parskip}
  }{% else
    \setlength{\parindent}{0pt}
    \setlength{\parskip}{6pt plus 2pt minus 1pt}}
}{% if KOMA class
  \KOMAoptions{parskip=half}}
\makeatother
\usepackage{xcolor}
\usepackage{color}
\usepackage{fancyvrb}
\newcommand{\VerbBar}{|}
\newcommand{\VERB}{\Verb[commandchars=\\\{\}]}
\DefineVerbatimEnvironment{Highlighting}{Verbatim}{commandchars=\\\{\}}
% Add ',fontsize=\small' for more characters per line
\newenvironment{Shaded}{}{}
\newcommand{\AlertTok}[1]{\textcolor[rgb]{1.00,0.00,0.00}{\textbf{#1}}}
\newcommand{\AnnotationTok}[1]{\textcolor[rgb]{0.38,0.63,0.69}{\textbf{\textit{#1}}}}
\newcommand{\AttributeTok}[1]{\textcolor[rgb]{0.49,0.56,0.16}{#1}}
\newcommand{\BaseNTok}[1]{\textcolor[rgb]{0.25,0.63,0.44}{#1}}
\newcommand{\BuiltInTok}[1]{\textcolor[rgb]{0.00,0.50,0.00}{#1}}
\newcommand{\CharTok}[1]{\textcolor[rgb]{0.25,0.44,0.63}{#1}}
\newcommand{\CommentTok}[1]{\textcolor[rgb]{0.38,0.63,0.69}{\textit{#1}}}
\newcommand{\CommentVarTok}[1]{\textcolor[rgb]{0.38,0.63,0.69}{\textbf{\textit{#1}}}}
\newcommand{\ConstantTok}[1]{\textcolor[rgb]{0.53,0.00,0.00}{#1}}
\newcommand{\ControlFlowTok}[1]{\textcolor[rgb]{0.00,0.44,0.13}{\textbf{#1}}}
\newcommand{\DataTypeTok}[1]{\textcolor[rgb]{0.56,0.13,0.00}{#1}}
\newcommand{\DecValTok}[1]{\textcolor[rgb]{0.25,0.63,0.44}{#1}}
\newcommand{\DocumentationTok}[1]{\textcolor[rgb]{0.73,0.13,0.13}{\textit{#1}}}
\newcommand{\ErrorTok}[1]{\textcolor[rgb]{1.00,0.00,0.00}{\textbf{#1}}}
\newcommand{\ExtensionTok}[1]{#1}
\newcommand{\FloatTok}[1]{\textcolor[rgb]{0.25,0.63,0.44}{#1}}
\newcommand{\FunctionTok}[1]{\textcolor[rgb]{0.02,0.16,0.49}{#1}}
\newcommand{\ImportTok}[1]{\textcolor[rgb]{0.00,0.50,0.00}{\textbf{#1}}}
\newcommand{\InformationTok}[1]{\textcolor[rgb]{0.38,0.63,0.69}{\textbf{\textit{#1}}}}
\newcommand{\KeywordTok}[1]{\textcolor[rgb]{0.00,0.44,0.13}{\textbf{#1}}}
\newcommand{\NormalTok}[1]{#1}
\newcommand{\OperatorTok}[1]{\textcolor[rgb]{0.40,0.40,0.40}{#1}}
\newcommand{\OtherTok}[1]{\textcolor[rgb]{0.00,0.44,0.13}{#1}}
\newcommand{\PreprocessorTok}[1]{\textcolor[rgb]{0.74,0.48,0.00}{#1}}
\newcommand{\RegionMarkerTok}[1]{#1}
\newcommand{\SpecialCharTok}[1]{\textcolor[rgb]{0.25,0.44,0.63}{#1}}
\newcommand{\SpecialStringTok}[1]{\textcolor[rgb]{0.73,0.40,0.53}{#1}}
\newcommand{\StringTok}[1]{\textcolor[rgb]{0.25,0.44,0.63}{#1}}
\newcommand{\VariableTok}[1]{\textcolor[rgb]{0.10,0.09,0.49}{#1}}
\newcommand{\VerbatimStringTok}[1]{\textcolor[rgb]{0.25,0.44,0.63}{#1}}
\newcommand{\WarningTok}[1]{\textcolor[rgb]{0.38,0.63,0.69}{\textbf{\textit{#1}}}}
\usepackage{longtable,booktabs,array}
\usepackage{calc} % for calculating minipage widths
% Correct order of tables after \paragraph or \subparagraph
\usepackage{etoolbox}
\makeatletter
\patchcmd\longtable{\par}{\if@noskipsec\mbox{}\fi\par}{}{}
\makeatother
% Allow footnotes in longtable head/foot
\IfFileExists{footnotehyper.sty}{\usepackage{footnotehyper}}{\usepackage{footnote}}
\makesavenoteenv{longtable}
\setlength{\emergencystretch}{3em} % prevent overfull lines
\providecommand{\tightlist}{%
  \setlength{\itemsep}{0pt}\setlength{\parskip}{0pt}}
\setcounter{secnumdepth}{-\maxdimen} % remove section numbering
\ifLuaTeX
  \usepackage{selnolig}  % disable illegal ligatures
\fi
\IfFileExists{bookmark.sty}{\usepackage{bookmark}}{\usepackage{hyperref}}
\IfFileExists{xurl.sty}{\usepackage{xurl}}{} % add URL line breaks if available
\urlstyle{same}
\hypersetup{
  hidelinks,
  pdfcreator={LaTeX via pandoc}}

\author{}
\date{}

\begin{document}

\hypertarget{satspath}{%
\section{SatsPath}\label{satspath}}

\hypertarget{funcionamiento-actual-capacidades-y-camino-hacia-producciuxf3n}{%
\subsection{Funcionamiento actual, capacidades y camino hacia
producción}\label{funcionamiento-actual-capacidades-y-camino-hacia-producciuxf3n}}

\textbf{Repositorio analizado:} \texttt{Truja503/satspath}\\
\textbf{Estado del análisis:} julio de 2026\\
\textbf{Objetivo del documento:} explicar qué es SatsPath, cómo funciona
actualmente, qué puede lograr con su implementación presente y qué
componentes faltan para convertirlo en infraestructura utilizada por
wallets y aplicaciones reales.

\begin{center}\rule{0.5\linewidth}{0.5pt}\end{center}

\hypertarget{resumen-ejecutivo}{%
\subsection{1. Resumen ejecutivo}\label{resumen-ejecutivo}}

SatsPath es un protocolo y conjunto de herramientas para descubrir,
verificar y seleccionar métodos de pago de Bitcoin a partir de un
identificador legible, por ejemplo:

\begin{Shaded}
\begin{Highlighting}[]
\NormalTok{rodrigo@satspath.dev}
\end{Highlighting}
\end{Shaded}

En lugar de obligar al pagador a conocer previamente si el receptor
utiliza Lightning, una dirección on-chain, Ark, BOLT12 u otro método,
SatsPath permite publicar un perfil firmado que contiene los métodos
públicos disponibles.

El flujo general es:

\begin{Shaded}
\begin{Highlighting}[]
\NormalTok{Identificador}
\NormalTok{    ↓}
\NormalTok{Cadena de resolvers}
\NormalTok{    ↓}
\NormalTok{Perfil de pago firmado}
\NormalTok{    ↓}
\NormalTok{Verificación de firma, identidad y expiración}
\NormalTok{    ↓}
\NormalTok{Selección del método de pago}
\NormalTok{    ↓}
\NormalTok{Invoice, URI, QR o handoff hacia una wallet}
\end{Highlighting}
\end{Shaded}

SatsPath puede utilizar información pública real de Bitcoin y generar
instrucciones que una wallet puede pagar. Por ejemplo, puede obtener una
invoice BOLT11 mediante LNURL, construir un URI BIP-21 para una
dirección on-chain o devolver un pointer de Ark.

SatsPath no necesita almacenar seeds ni claves privadas de fondos para
cumplir su función. Las claves de identidad del protocolo sirven para
firmar perfiles públicos y demostrar que estos no fueron modificados.
Las claves que gastan Bitcoin permanecen dentro de la wallet del
usuario.

La propuesta central puede resumirse así:

\begin{quote}
SatsPath transforma una identidad legible en una instrucción de pago
verificable y selecciona el rail más apropiado para que una wallet
complete el pago.
\end{quote}

\begin{center}\rule{0.5\linewidth}{0.5pt}\end{center}

\hypertarget{problema-que-resuelve}{%
\subsection{2. Problema que resuelve}\label{problema-que-resuelve}}

Bitcoin ya no posee una sola forma de recibir pagos. Un usuario puede
aceptar fondos mediante:

\begin{itemize}
\tightlist
\item
  Lightning Address.
\item
  LNURL-pay.
\item
  BOLT11.
\item
  BOLT12.
\item
  Dirección Bitcoin on-chain.
\item
  Silent Payments.
\item
  Ark.
\item
  Otros sistemas compatibles o futuros adaptadores.
\end{itemize}

Cada método tiene diferencias en:

\begin{itemize}
\tightlist
\item
  Velocidad.
\item
  Comisión.
\item
  Privacidad.
\item
  Disponibilidad.
\item
  Liquidez.
\item
  Interoperabilidad.
\item
  Dependencias de infraestructura.
\end{itemize}

Actualmente, la mayoría de wallets obliga al usuario o desarrollador a
conocer de antemano cuál método utilizar. SatsPath introduce una capa
intermedia que responde tres preguntas:

\begin{enumerate}
\def\labelenumi{\arabic{enumi}.}
\tightlist
\item
  ¿Quién es el receptor?
\item
  ¿Qué métodos públicos acepta?
\item
  ¿Cuál de esos métodos debería utilizarse para este pago?
\end{enumerate}

\begin{center}\rule{0.5\linewidth}{0.5pt}\end{center}

\hypertarget{quuxe9-es-un-perfil-de-pago}{%
\subsection{3. Qué es un perfil de
pago}\label{quuxe9-es-un-perfil-de-pago}}

El elemento principal del protocolo es el \texttt{SignedPaymentProfile}.

Un perfil puede incluir:

\begin{Shaded}
\begin{Highlighting}[]
\FunctionTok{\{}
  \DataTypeTok{"alias"}\FunctionTok{:} \StringTok{"rodrigo@satspath.dev"}\FunctionTok{,}
  \DataTypeTok{"identity\_pubkey"}\FunctionTok{:} \StringTok{"02..."}\FunctionTok{,}
  \DataTypeTok{"methods"}\FunctionTok{:} \OtherTok{[}
    \FunctionTok{\{}
      \DataTypeTok{"type"}\FunctionTok{:} \StringTok{"Lightning"}\FunctionTok{,}
      \DataTypeTok{"lightning\_address"}\FunctionTok{:} \StringTok{"rodrigo@satspath.dev"}
    \FunctionTok{\}}\OtherTok{,}
    \FunctionTok{\{}
      \DataTypeTok{"type"}\FunctionTok{:} \StringTok{"Onchain"}\FunctionTok{,}
      \DataTypeTok{"network"}\FunctionTok{:} \StringTok{"Mainnet"}\FunctionTok{,}
      \DataTypeTok{"address"}\FunctionTok{:} \StringTok{"bc1q..."}
    \FunctionTok{\}}\OtherTok{,}
    \FunctionTok{\{}
      \DataTypeTok{"type"}\FunctionTok{:} \StringTok{"Ark"}\FunctionTok{,}
      \DataTypeTok{"server"}\FunctionTok{:} \StringTok{"https://ark.example.com"}\FunctionTok{,}
      \DataTypeTok{"pubkey"}\FunctionTok{:} \StringTok{"02..."}
    \FunctionTok{\}}
  \OtherTok{]}\FunctionTok{,}
  \DataTypeTok{"updated\_at"}\FunctionTok{:} \DecValTok{1782810000}\FunctionTok{,}
  \DataTypeTok{"expires\_at"}\FunctionTok{:} \DecValTok{1785402000}\FunctionTok{,}
  \DataTypeTok{"sequence"}\FunctionTok{:} \DecValTok{4}
\FunctionTok{\}}
\end{Highlighting}
\end{Shaded}

El perfil se firma con una clave de identidad secp256k1.

\begin{Shaded}
\begin{Highlighting}[]
\NormalTok{profile}
\NormalTok{    ↓ canonical JSON}
\NormalTok{SHA{-}256}
\NormalTok{    ↓}
\NormalTok{Firma ECDSA secp256k1}
\end{Highlighting}
\end{Shaded}

La firma permite comprobar que:

\begin{itemize}
\tightlist
\item
  El perfil no fue modificado.
\item
  Los métodos no fueron reemplazados durante la resolución.
\item
  La clave pública de identidad corresponde con la firma.
\item
  La versión recibida conserva integridad criptográfica.
\end{itemize}

La clave de identidad de SatsPath no tiene que controlar fondos. Su
función es firmar información pública, no firmar transacciones de
Bitcoin.

Implementación principal:

\begin{Shaded}
\begin{Highlighting}[]
\NormalTok{crates/satspath{-}core/src/profile.rs}
\NormalTok{crates/satspath{-}core/src/crypto.rs}
\NormalTok{crates/satspath{-}core/src/validation.rs}
\end{Highlighting}
\end{Shaded}

\begin{center}\rule{0.5\linewidth}{0.5pt}\end{center}

\hypertarget{funcionamiento-completo}{%
\subsection{4. Funcionamiento completo}\label{funcionamiento-completo}}

\hypertarget{creaciuxf3n-de-identidad}{%
\subsubsection{4.1 Creación de
identidad}\label{creaciuxf3n-de-identidad}}

El usuario genera localmente una clave de identidad para SatsPath.

Esta clave permite:

\begin{itemize}
\tightlist
\item
  Firmar perfiles.
\item
  Firmar invitaciones.
\item
  Identificar actualizaciones legítimas.
\item
  Rotar la identidad en versiones posteriores.
\end{itemize}

No sustituye las claves de una wallet ni debe utilizarse para controlar
UTXOs o canales Lightning.

\hypertarget{registro-de-muxe9todos-de-recepciuxf3n}{%
\subsubsection{4.2 Registro de métodos de
recepción}\label{registro-de-muxe9todos-de-recepciuxf3n}}

El receptor agrega uno o varios métodos públicos:

\begin{Shaded}
\begin{Highlighting}[]
\NormalTok{Lightning Address}
\NormalTok{Dirección on{-}chain}
\NormalTok{Silent Payment address o pubkey}
\NormalTok{BOLT12 offer}
\NormalTok{Ark server + pubkey o receive pointer}
\end{Highlighting}
\end{Shaded}

SatsPath valida el formato y construye un perfil público.

\hypertarget{firma-del-perfil}{%
\subsubsection{4.3 Firma del perfil}\label{firma-del-perfil}}

El perfil se serializa y firma con la clave de identidad.

El resultado es un objeto que puede publicarse por varios medios sin
confiar plenamente en el transporte.

\hypertarget{publicaciuxf3n}{%
\subsubsection{4.4 Publicación}\label{publicaciuxf3n}}

El perfil firmado puede distribuirse mediante:

\begin{itemize}
\tightlist
\item
  Registro local.
\item
  Endpoint HTTPS \texttt{.well-known}.
\item
  BIP-353 mediante DNS.
\item
  Nostr y NIP-05.
\item
  Pear/Holepunch P2P.
\item
  Exportación e importación manual.
\item
  Un futuro resolver de plataforma.
\end{itemize}

La arquitectura es neutral al transporte. HTTPS, DNS, Nostr y P2P son
formas de mover el perfil, pero la confianza final proviene de la firma
y de las pruebas adicionales de identidad.

\hypertarget{resoluciuxf3n}{%
\subsubsection{4.5 Resolución}\label{resoluciuxf3n}}

El pagador introduce un identificador:

\begin{Shaded}
\begin{Highlighting}[]
\NormalTok{alice@example.com}
\end{Highlighting}
\end{Shaded}

SatsPath consulta una cadena de resolvers hasta encontrar un perfil.

La cadena actual puede incluir:

\begin{Shaded}
\begin{Highlighting}[]
\NormalTok{Registro local}
\NormalTok{    ↓}
\NormalTok{BIP{-}353}
\NormalTok{    ↓}
\NormalTok{HTTPS well{-}known}
\NormalTok{    ↓}
\NormalTok{Nostr / NIP{-}05}
\NormalTok{    ↓}
\NormalTok{Pear / Holepunch P2P}
\end{Highlighting}
\end{Shaded}

El resolver también compara el alias solicitado con el alias contenido
en el perfil para evitar ataques de sustitución.

Implementación principal:

\begin{Shaded}
\begin{Highlighting}[]
\NormalTok{crates/satspath{-}core/src/resolver.rs}
\NormalTok{crates/satspath{-}core/src/resolvers/}
\NormalTok{crates/satspath{-}core/src/registry.rs}
\NormalTok{crates/satspath{-}core/src/peer\_registry.rs}
\end{Highlighting}
\end{Shaded}

\hypertarget{verificaciuxf3n}{%
\subsubsection{4.6 Verificación}\label{verificaciuxf3n}}

Antes de utilizar el perfil, SatsPath puede comprobar:

\begin{itemize}
\tightlist
\item
  Firma criptográfica.
\item
  Alias correcto.
\item
  Expiración.
\item
  Número de secuencia.
\item
  Presencia de material privado prohibido.
\item
  Pruebas de propiedad de métodos.
\item
  Procedencia del resolver cuando esté disponible.
\item
  Compatibilidad de red de una dirección Bitcoin.
\end{itemize}

Una firma válida demuestra control de la identidad de protocolo, pero no
necesariamente control del correo electrónico o dominio utilizado en el
alias. Esa segunda capa requiere DNSSEC, NIP-05, HTTPS verificado o un
challenge de plataforma.

\hypertarget{selecciuxf3n-de-rail}{%
\subsubsection{4.7 Selección de rail}\label{selecciuxf3n-de-rail}}

Después de resolver y verificar el perfil, el router analiza los métodos
disponibles.

Actualmente existen políticas para:

\begin{itemize}
\tightlist
\item
  Seleccionar Lightning en pagos pequeños cuando está disponible.
\item
  Utilizar on-chain según las comisiones estimadas.
\item
  Utilizar Ark como alternativa.
\item
  Generar razones legibles para la decisión.
\item
  Evaluar urgencia y snapshots de fees.
\item
  Preparar rutas divididas en módulos experimentales.
\end{itemize}

Implementación principal:

\begin{Shaded}
\begin{Highlighting}[]
\NormalTok{crates/satspath{-}router/src/router.rs}
\NormalTok{crates/satspath{-}router/src/scoring.rs}
\NormalTok{crates/satspath{-}router/src/priority.rs}
\NormalTok{crates/satspath{-}router/src/fees.rs}
\end{Highlighting}
\end{Shaded}

\hypertarget{construcciuxf3n-de-la-instrucciuxf3n-de-pago}{%
\subsubsection{4.8 Construcción de la instrucción de
pago}\label{construcciuxf3n-de-la-instrucciuxf3n-de-pago}}

Dependiendo del método elegido, SatsPath puede producir:

\hypertarget{lightning}{%
\paragraph{Lightning}\label{lightning}}

\begin{itemize}
\tightlist
\item
  Lightning Address.
\item
  LNURL.
\item
  Invoice BOLT11 obtenida desde un callback LNURL-pay.
\item
  BOLT12 offer o información relacionada.
\end{itemize}

\hypertarget{on-chain}{%
\paragraph{On-chain}\label{on-chain}}

\begin{Shaded}
\begin{Highlighting}[]
\NormalTok{bitcoin:bc1q...?amount=0.00021000}
\end{Highlighting}
\end{Shaded}

\hypertarget{ark}{%
\paragraph{Ark}\label{ark}}

\begin{Shaded}
\begin{Highlighting}[]
\NormalTok{ark:\textless{}pubkey\textgreater{}?server=\textless{}server\textgreater{}\&amount=\textless{}sats\textgreater{}}
\end{Highlighting}
\end{Shaded}

También puede generar un QR SVG y una respuesta JSON estable para el
frontend.

\hypertarget{pago-real-mediante-una-wallet}{%
\subsubsection{4.9 Pago real mediante una
wallet}\label{pago-real-mediante-una-wallet}}

La instrucción producida por SatsPath puede ser entregada a una wallet
que controle fondos.

\begin{Shaded}
\begin{Highlighting}[]
\NormalTok{SatsPath}
\NormalTok{    ↓ invoice / URI / pointer}
\NormalTok{Wallet del pagador}
\NormalTok{    ↓ firma y ejecución}
\NormalTok{Bitcoin, Lightning o Ark}
\end{Highlighting}
\end{Shaded}

Esto permite realizar pagos reales sin que SatsPath almacene las claves
privadas de los fondos.

La responsabilidad se divide así:

\begin{longtable}[]{@{}
  >{\raggedright\arraybackslash}p{(\columnwidth - 2\tabcolsep) * \real{0.5000}}
  >{\raggedright\arraybackslash}p{(\columnwidth - 2\tabcolsep) * \real{0.5000}}@{}}
\toprule\noalign{}
\begin{minipage}[b]{\linewidth}\raggedright
Componente
\end{minipage} & \begin{minipage}[b]{\linewidth}\raggedright
Responsabilidad
\end{minipage} \\
\midrule\noalign{}
\endhead
\bottomrule\noalign{}
\endlastfoot
SatsPath & Resolver identidad, verificar perfil, escoger rail y
construir la instrucción \\
Wallet & Mostrar confirmación, firmar y ejecutar el pago \\
Red o protocolo & Liquidar o transmitir el pago \\
\end{longtable}

Esta separación es una propiedad del diseño, no una limitación
conceptual. Mantiene a SatsPath como infraestructura no custodial y
facilita su integración con varias wallets.

\begin{center}\rule{0.5\linewidth}{0.5pt}\end{center}

\hypertarget{quuxe9-existe-actualmente}{%
\subsection{5. Qué existe actualmente}\label{quuxe9-existe-actualmente}}

\hypertarget{core-del-protocolo}{%
\subsubsection{5.1 Core del protocolo}\label{core-del-protocolo}}

Existe implementación para:

\begin{itemize}
\tightlist
\item
  \texttt{PaymentProfile}.
\item
  \texttt{SignedPaymentProfile}.
\item
  Métodos Lightning, on-chain y Ark.
\item
  Firma y verificación secp256k1.
\item
  Validación de datos públicos.
\item
  Invitaciones.
\item
  Expiración.
\item
  Nonces.
\item
  Secuencias.
\item
  Pruebas de propiedad.
\item
  Rotación de claves.
\item
  Codificación de requests universales.
\end{itemize}

\hypertarget{resolver-chain}{%
\subsubsection{5.2 Resolver chain}\label{resolver-chain}}

Existen resolvers para:

\begin{itemize}
\tightlist
\item
  Registro local.
\item
  HTTPS.
\item
  BIP-353.
\item
  Nostr y NIP-05.
\item
  P2P mediante el SDK Pear/Holepunch en el daemon.
\end{itemize}

El resolver HTTPS verifica firma y expiración antes de aceptar un
perfil.

\hypertarget{router}{%
\subsubsection{5.3 Router}\label{router}}

El router puede:

\begin{itemize}
\tightlist
\item
  Detectar métodos disponibles.
\item
  Consultar fees on-chain.
\item
  Elegir un rail.
\item
  Calcular una comisión aproximada.
\item
  Devolver ETA y explicación.
\item
  Construir un payload de pago.
\item
  Producir un contrato \texttt{QuoteResponse} estable.
\end{itemize}

Estados actuales:

\begin{Shaded}
\begin{Highlighting}[]
\NormalTok{ok}
\NormalTok{not\_registered}
\NormalTok{no\_route}
\NormalTok{invalid\_signature}
\end{Highlighting}
\end{Shaded}

\hypertarget{lnurl-y-bolt11}{%
\subsubsection{5.4 LNURL y BOLT11}\label{lnurl-y-bolt11}}

La implementación puede:

\begin{enumerate}
\def\labelenumi{\arabic{enumi}.}
\tightlist
\item
  Resolver una Lightning Address.
\item
  Obtener metadata LNURL-pay.
\item
  Validar \texttt{minSendable} y \texttt{maxSendable}.
\item
  Solicitar una invoice.
\item
  Parsear BOLT11.
\item
  Verificar el monto.
\item
  Verificar expiración.
\item
  Entregar la invoice como payload o QR.
\end{enumerate}

\hypertarget{daemon-local}{%
\subsubsection{5.5 Daemon local}\label{daemon-local}}

\texttt{satspathd} expone una API HTTP con endpoints para:

\begin{Shaded}
\begin{Highlighting}[]
\NormalTok{GET  /health}
\NormalTok{GET  /v1/status}
\NormalTok{GET  /v1/node}
\NormalTok{GET  /v1/profile}
\NormalTok{POST /v1/profile}
\NormalTok{POST /v1/profile/methods}
\NormalTok{POST /v1/resolve}
\NormalTok{POST /v1/quote}
\NormalTok{POST /v1/pay}
\NormalTok{POST /v1/send}
\NormalTok{POST /v1/receive}
\NormalTok{POST /v1/dns/resolve}
\NormalTok{GET  /v1/peers}
\NormalTok{GET  /v1/connections}
\end{Highlighting}
\end{Shaded}

El daemon también incluye una UI local, generación de QR, administración
de perfiles y un bridge hacia el SDK P2P.

\hypertarget{sdks-y-bindings}{%
\subsubsection{5.6 SDKs y bindings}\label{sdks-y-bindings}}

El workspace contiene:

\begin{Shaded}
\begin{Highlighting}[]
\NormalTok{satspath{-}core}
\NormalTok{satspath{-}router}
\NormalTok{satspath{-}cli}
\NormalTok{satspath{-}swaps}
\NormalTok{satspathd}
\NormalTok{satspath{-}wasm}
\NormalTok{satspath{-}ffi}
\end{Highlighting}
\end{Shaded}

También existen bindings generados para Kotlin y Swift mediante UniFFI,
además de funciones WASM para aplicaciones web.

\hypertarget{p2p}{%
\subsubsection{5.7 P2P}\label{p2p}}

El SDK de Pear/Holepunch permite publicar y resolver perfiles firmados
mediante Hyperswarm.

El daemon puede iniciar procesos Node para:

\begin{itemize}
\tightlist
\item
  Publicar un perfil.
\item
  Resolver perfiles.
\item
  Validar el perfil nuevamente en Rust.
\item
  Guardar perfiles resueltos en registros locales.
\end{itemize}

\hypertarget{funciones-avanzadas-presentes}{%
\subsubsection{5.8 Funciones avanzadas
presentes}\label{funciones-avanzadas-presentes}}

El repositorio contiene módulos para:

\begin{itemize}
\tightlist
\item
  BOLT12.
\item
  Silent Payments.
\item
  Split Payments.
\item
  Key Rotation.
\item
  Ownership proofs.
\item
  Ark intents y rutas.
\item
  WASM.
\item
  FFI móvil.
\end{itemize}

Algunas de estas funciones están implementadas a nivel de tipos, parsing
o planificación, pero aún necesitan interoperabilidad completa con
wallets reales.

\begin{center}\rule{0.5\linewidth}{0.5pt}\end{center}

\hypertarget{quuxe9-puede-lograr-satspath}{%
\subsection{6. Qué puede lograr
SatsPath}\label{quuxe9-puede-lograr-satspath}}

\hypertarget{identidad-de-pago-universal}{%
\subsubsection{6.1 Identidad de pago
universal}\label{identidad-de-pago-universal}}

Un usuario podría compartir una sola identidad:

\begin{Shaded}
\begin{Highlighting}[]
\NormalTok{rodrigo@satspath.dev}
\end{Highlighting}
\end{Shaded}

sin publicar manualmente una dirección diferente para cada protocolo.

\hypertarget{descubrimiento-automuxe1tico}{%
\subsubsection{6.2 Descubrimiento
automático}\label{descubrimiento-automuxe1tico}}

Una wallet podría descubrir qué métodos acepta el receptor sin pedirle
que seleccione manualmente Lightning, Ark u on-chain.

\hypertarget{selecciuxf3n-inteligente}{%
\subsubsection{6.3 Selección
inteligente}\label{selecciuxf3n-inteligente}}

El router puede evolucionar para elegir según:

\begin{itemize}
\tightlist
\item
  Fee máximo.
\item
  Urgencia.
\item
  Privacidad.
\item
  Liquidez Lightning.
\item
  Salud del rail.
\item
  Monto.
\item
  Preferencias del usuario.
\item
  Compatibilidad de la wallet.
\item
  Confianza del resolver.
\end{itemize}

\hypertarget{interoperabilidad-entre-wallets}{%
\subsubsection{6.4 Interoperabilidad entre
wallets}\label{interoperabilidad-entre-wallets}}

SatsPath puede convertirse en un SDK que varias wallets integren para
obtener el mismo resultado a partir del mismo perfil firmado.

\hypertarget{comercios}{%
\subsubsection{6.5 Comercios}\label{comercios}}

Un comercio podría publicar un perfil con múltiples métodos y permitir
que cada cliente pague mediante la ruta compatible con su wallet.

\hypertarget{fallback-entre-rails}{%
\subsubsection{6.6 Fallback entre rails}\label{fallback-entre-rails}}

Si Lightning no está disponible, la wallet podría utilizar Ark u
on-chain sin solicitar otra dirección al receptor.

\hypertarget{infraestructura-no-custodial}{%
\subsubsection{6.7 Infraestructura no
custodial}\label{infraestructura-no-custodial}}

SatsPath puede operar sin controlar fondos, reduciendo riesgos
regulatorios y de seguridad asociados con custodiar Bitcoin.

\hypertarget{base-para-protocolos-adicionales}{%
\subsubsection{6.8 Base para protocolos
adicionales}\label{base-para-protocolos-adicionales}}

El sistema puede agregar adaptadores para nuevos rails siempre que estos
puedan representarse mediante instrucciones públicas verificables.

\begin{center}\rule{0.5\linewidth}{0.5pt}\end{center}

\hypertarget{quuxe9-hace-falta}{%
\subsection{7. Qué hace falta}\label{quuxe9-hace-falta}}

\hypertarget{verificaciuxf3n-real-de-identidad}{%
\subsection{7.1 Verificación real de
identidad}\label{verificaciuxf3n-real-de-identidad}}

El daemon actual utiliza un verificador simulado para challenges de
correo.

Hace falta implementar:

\begin{itemize}
\tightlist
\item
  Magic links reales.
\item
  Tokens aleatorios de un solo uso.
\item
  Hash del token en almacenamiento.
\item
  Expiración.
\item
  Rate limiting.
\item
  Protección contra enumeración.
\item
  Confirmación explícita del alias.
\item
  Pruebas de control de dominio.
\end{itemize}

Para dominios propios se debe priorizar:

\begin{itemize}
\tightlist
\item
  DNSSEC y BIP-353.
\item
  HTTPS \texttt{.well-known}.
\item
  NIP-05.
\end{itemize}

Para correos de consumo se necesita un resolver de plataforma que valide
la cuenta antes de publicar el perfil.

\hypertarget{procedencia-y-nivel-de-confianza}{%
\subsection{7.2 Procedencia y nivel de
confianza}\label{procedencia-y-nivel-de-confianza}}

El resolver debería devolver no solamente el perfil, sino también su
procedencia.

Propuesta:

\begin{Shaded}
\begin{Highlighting}[]
\KeywordTok{struct}\NormalTok{ ResolvedProfile }\OperatorTok{\{}
\NormalTok{    profile}\OperatorTok{:}\NormalTok{ SignedPaymentProfile}\OperatorTok{,}
\NormalTok{    source}\OperatorTok{:}\NormalTok{ ResolverSource}\OperatorTok{,}
\NormalTok{    identifier\_verified}\OperatorTok{:} \DataTypeTok{bool}\OperatorTok{,}
\NormalTok{    verification\_method}\OperatorTok{:}\NormalTok{ VerificationMethod}\OperatorTok{,}
\NormalTok{    trust\_level}\OperatorTok{:}\NormalTok{ TrustLevel}\OperatorTok{,}
\NormalTok{    fetched\_at}\OperatorTok{:} \DataTypeTok{i64}\OperatorTok{,}
\OperatorTok{\}}
\end{Highlighting}
\end{Shaded}

Posibles fuentes:

\begin{Shaded}
\begin{Highlighting}[]
\NormalTok{local}
\NormalTok{https}
\NormalTok{bip353\_dnssec}
\NormalTok{nostr\_nip05}
\NormalTok{p2p}
\NormalTok{platform}
\end{Highlighting}
\end{Shaded}

La aplicación necesita distinguir entre:

\begin{itemize}
\tightlist
\item
  Perfil firmado.
\item
  Dominio verificado.
\item
  Correo verificado.
\item
  Método de pago verificado.
\item
  Perfil recibido desde una fuente no autenticada.
\end{itemize}

\hypertarget{estados-de-error-muxe1s-precisos}{%
\subsection{7.3 Estados de error más
precisos}\label{estados-de-error-muxe1s-precisos}}

Actualmente varios errores de resolución pueden terminar representados
como \texttt{not\_registered}.

Se recomienda ampliar el contrato:

\begin{Shaded}
\begin{Highlighting}[]
\NormalTok{ok}
\NormalTok{not\_found}
\NormalTok{temporarily\_unavailable}
\NormalTok{invalid\_signature}
\NormalTok{expired\_profile}
\NormalTok{identifier\_unverified}
\NormalTok{unsupported\_method}
\NormalTok{no\_route}
\NormalTok{resolver\_error}
\end{Highlighting}
\end{Shaded}

Esto evitará confundir una caída temporal de red con un usuario que no
existe.

\hypertarget{seguridad-del-daemon}{%
\subsection{7.4 Seguridad del daemon}\label{seguridad-del-daemon}}

Antes de exponer \texttt{satspathd} fuera de localhost hacen falta:

\begin{itemize}
\tightlist
\item
  Autenticación de administración.
\item
  Token de sesión local.
\item
  Protección CSRF.
\item
  CORS restringido.
\item
  Límites de requests.
\item
  Límite de tamaño de payload.
\item
  Separación entre API pública y API administrativa.
\item
  Restricción explícita para binds públicos.
\item
  Logs sin datos sensibles.
\item
  Manejo seguro de procesos P2P secundarios.
\end{itemize}

\hypertarget{protecciuxf3n-de-la-clave-de-identidad}{%
\subsection{7.5 Protección de la clave de
identidad}\label{protecciuxf3n-de-la-clave-de-identidad}}

La clave de identidad no controla fondos, pero sí puede firmar perfiles
falsos si es robada.

Se recomienda integrar:

\begin{itemize}
\tightlist
\item
  macOS Keychain.
\item
  iOS Keychain.
\item
  Android Keystore.
\item
  Linux Secret Service.
\item
  Windows Credential Manager.
\item
  Cifrado AES-GCM como fallback.
\end{itemize}

También hace falta:

\begin{itemize}
\tightlist
\item
  Backup.
\item
  Recuperación.
\item
  Revocación.
\item
  Rotación asistida.
\item
  Confirmación en un segundo dispositivo.
\end{itemize}

\hypertarget{router-basado-en-capacidades-reales}{%
\subsection{7.6 Router basado en capacidades
reales}\label{router-basado-en-capacidades-reales}}

La política actual utiliza umbrales generales. Para producción debe
evolucionar hacia un sistema de scoring configurable.

Factores recomendados:

\begin{Shaded}
\begin{Highlighting}[]
\NormalTok{Costo estimado}
\NormalTok{Probabilidad de éxito}
\NormalTok{Tiempo esperado}
\NormalTok{Privacidad}
\NormalTok{Monto mínimo y máximo}
\NormalTok{Liquidez disponible}
\NormalTok{Disponibilidad del servidor}
\NormalTok{Confianza del método}
\NormalTok{Compatibilidad de la wallet}
\NormalTok{Urgencia}
\NormalTok{Preferencias del usuario}
\end{Highlighting}
\end{Shaded}

El router también debe respetar realmente:

\begin{Shaded}
\begin{Highlighting}[]
\NormalTok{max\_fee\_sats}
\NormalTok{max\_fee\_percent}
\NormalTok{preferred\_rails}
\NormalTok{blocked\_rails}
\NormalTok{privacy\_mode}
\end{Highlighting}
\end{Shaded}

\hypertarget{disponibilidad-real-de-lightning}{%
\subsection{7.7 Disponibilidad real de
Lightning}\label{disponibilidad-real-de-lightning}}

Hace falta reemplazar varias suposiciones estáticas por comprobaciones
reales:

\begin{itemize}
\tightlist
\item
  Límites de LNURL.
\item
  Capacidad de obtener invoice.
\item
  Validez del callback.
\item
  Expiración.
\item
  Compatibilidad de red.
\item
  Soporte de BOLT12.
\item
  Resultado del intento de pago reportado por la wallet.
\end{itemize}

SatsPath no tiene que inspeccionar canales privados, pero la wallet
integradora puede aportar señales sobre probabilidad de éxito y
liquidez.

\hypertarget{protecciuxf3n-de-red-y-ssrf}{%
\subsection{7.8 Protección de red y
SSRF}\label{protecciuxf3n-de-red-y-ssrf}}

Los resolvers HTTP, Nostr y LNURL deben aplicar políticas estrictas:

\begin{itemize}
\tightlist
\item
  HTTPS obligatorio salvo desarrollo local.
\item
  Bloqueo de localhost y redes privadas para contenido remoto.
\item
  Control de redirects.
\item
  Timeouts.
\item
  Límite de tamaño de respuesta.
\item
  Validación del content type.
\item
  Allowlist o reglas de dominio para callbacks.
\item
  Resolución DNS segura.
\item
  Protección contra rebinding.
\end{itemize}

\hypertarget{bip-353-y-dnssec-operativos}{%
\subsection{7.9 BIP-353 y DNSSEC
operativos}\label{bip-353-y-dnssec-operativos}}

Hace falta convertir el soporte actual en una experiencia operativa
completa:

\begin{itemize}
\tightlist
\item
  Resolver DNSSEC confiable.
\item
  Publicación automatizada.
\item
  Compatibilidad con proveedores DNS.
\item
  Manejo de TTL.
\item
  Rotación de registros.
\item
  Validación de redes.
\item
  Herramienta de diagnóstico.
\end{itemize}

\hypertarget{publicaciuxf3n-nostr-desde-las-libreruxedas-principales}{%
\subsection{7.10 Publicación Nostr desde las librerías
principales}\label{publicaciuxf3n-nostr-desde-las-libreruxedas-principales}}

El resolver Nostr ya puede consumir perfiles, pero la publicación
necesita una interfaz integrada y coherente para:

\begin{itemize}
\tightlist
\item
  Crear evento.
\item
  Firmarlo.
\item
  Publicar en varios relays.
\item
  Confirmar recepción.
\item
  Actualizar perfiles.
\item
  Eliminar o invalidar versiones antiguas.
\end{itemize}

\hypertarget{p2p-embebido}{%
\subsection{7.11 P2P embebido}\label{p2p-embebido}}

El daemon posee un bridge funcional mediante procesos Node, pero el
binding FFI todavía mantiene un placeholder para P2P embebido.

Hace falta decidir una dirección:

\begin{enumerate}
\def\labelenumi{\arabic{enumi}.}
\tightlist
\item
  Mantener Hyperswarm como sidecar.
\item
  Integrar un runtime embebido.
\item
  Implementar un \texttt{P2pTransport} abstracto con múltiples backends.
\end{enumerate}

La API debe evitar depender de rutas relativas del repositorio o scripts
locales cuando se distribuya como aplicación instalada.

\hypertarget{interoperabilidad-ark}{%
\subsection{7.12 Interoperabilidad Ark}\label{interoperabilidad-ark}}

Para considerar Ark una integración completa hacen falta pruebas con una
wallet y servidor concretos.

Se debe comprobar:

\begin{itemize}
\tightlist
\item
  Formato del receive pointer.
\item
  URI aceptado.
\item
  QR compatible.
\item
  Expiración.
\item
  Servidor correcto.
\item
  Pubkey correcta.
\item
  Flujo Ark-to-Ark.
\item
  Fallback hacia Lightning u on-chain.
\item
  Manejo de errores.
\end{itemize}

No se deben inventar APIs de Arkade ni asumir formatos sin una prueba de
interoperabilidad.

\hypertarget{bolt12-completo}{%
\subsection{7.13 BOLT12 completo}\label{bolt12-completo}}

Hace falta validar el flujo completo:

\begin{Shaded}
\begin{Highlighting}[]
\NormalTok{Offer}
\NormalTok{    ↓}
\NormalTok{Invoice request}
\NormalTok{    ↓}
\NormalTok{Invoice}
\NormalTok{    ↓}
\NormalTok{Verificación}
\NormalTok{    ↓}
\NormalTok{Handoff a wallet}
\end{Highlighting}
\end{Shaded}

El parser por sí solo no es suficiente. Se necesitan test vectors y
pruebas con implementaciones Lightning existentes.

\hypertarget{silent-payments}{%
\subsection{7.14 Silent Payments}\label{silent-payments}}

Para integración completa hacen falta:

\begin{itemize}
\tightlist
\item
  Generación y validación compatible con BIP-352.
\item
  Network checks.
\item
  Payment URI correcto.
\item
  Compatibilidad con una wallet que pueda pagar.
\item
  Manejo de scan keys sin exponer material privado.
\item
  Test vectors oficiales.
\end{itemize}

\hypertarget{split-payments}{%
\subsection{7.15 Split Payments}\label{split-payments}}

El módulo debe definir claramente si divide:

\begin{itemize}
\tightlist
\item
  Un monto entre varios receptores.
\item
  Un pago entre varios rails.
\item
  Ingresos mediante políticas de porcentaje.
\item
  Pagos secuenciales con tolerancia parcial a fallos.
\end{itemize}

También hacen falta:

\begin{itemize}
\tightlist
\item
  Atomicidad o definición explícita de no atomicidad.
\item
  Idempotencia.
\item
  Reintentos.
\item
  Estado por receptor.
\item
  Política de redondeo.
\item
  Recuperación ante pagos parcialmente completados.
\end{itemize}

\hypertarget{revocaciuxf3n-y-rotaciuxf3n}{%
\subsection{7.16 Revocación y
rotación}\label{revocaciuxf3n-y-rotaciuxf3n}}

La rotación de claves necesita integrarse con los resolvers.

Un cliente debe poder saber:

\begin{itemize}
\tightlist
\item
  Que una clave anterior autorizó la nueva.
\item
  Cuál es la secuencia más reciente.
\item
  Si una clave fue revocada.
\item
  Si un perfil es un replay antiguo.
\item
  Qué fuente confirma la transición.
\end{itemize}

\hypertarget{cachuxe9-y-replay-protection}{%
\subsection{7.17 Caché y replay
protection}\label{cachuxe9-y-replay-protection}}

Cada perfil debería utilizar:

\begin{itemize}
\tightlist
\item
  \texttt{sequence} monotónica.
\item
  \texttt{updated\_at}.
\item
  \texttt{expires\_at}.
\item
  Nonce.
\item
  Clave efectiva después de rotación.
\end{itemize}

Los resolvers y caches deben rechazar versiones más antiguas que una
versión previamente observada.

\hypertarget{observabilidad}{%
\subsection{7.18 Observabilidad}\label{observabilidad}}

Hace falta instrumentar:

\begin{itemize}
\tightlist
\item
  Tiempo de resolución.
\item
  Resolver elegido.
\item
  Resolver fallido.
\item
  Perfil inválido.
\item
  Rail elegido.
\item
  Tiempo para obtener invoice.
\item
  Fallo de handoff.
\item
  Tasa de fallback.
\end{itemize}

Las métricas deben evitar almacenar aliases, direcciones o invoices
completas.

\hypertarget{ci-y-poluxedtica-de-releases}{%
\subsection{7.19 CI y política de
releases}\label{ci-y-poluxedtica-de-releases}}

El repositorio ya contiene un workflow de CI, pero debe convertirse en
un gate obligatorio.

Requisitos recomendados:

\begin{Shaded}
\begin{Highlighting}[]
\NormalTok{cargo fmt {-}{-}all {-}{-} {-}{-}check}
\NormalTok{cargo build {-}{-}workspace}
\NormalTok{cargo test {-}{-}workspace {-}{-}all{-}targets}
\NormalTok{cargo clippy {-}{-}workspace {-}{-}all{-}targets {-}{-} {-}D warnings}
\NormalTok{cargo audit}
\NormalTok{cargo deny}
\NormalTok{Tests WASM}
\NormalTok{Tests FFI}
\NormalTok{Tests del SDK JavaScript}
\end{Highlighting}
\end{Shaded}

También hacen falta:

\begin{itemize}
\tightlist
\item
  Branch protection.
\item
  Reviews obligatorias.
\item
  PRs pequeños.
\item
  Checklist con evidencia.
\item
  Releases versionadas.
\item
  Changelog.
\item
  Artefactos firmados.
\item
  Builds reproducibles.
\end{itemize}

\hypertarget{auditoruxeda-y-fuzzing}{%
\subsection{7.20 Auditoría y fuzzing}\label{auditoruxeda-y-fuzzing}}

Antes de una adopción amplia se recomienda:

\begin{itemize}
\tightlist
\item
  Auditoría externa de criptografía y parsing.
\item
  Fuzzing de perfiles.
\item
  Fuzzing de URIs.
\item
  Fuzzing de respuestas HTTP y Nostr.
\item
  Fuzzing de BOLT11 y BOLT12.
\item
  Pruebas de perfiles gigantes.
\item
  Pruebas de Unicode y normalización.
\item
  Pruebas de colisión y canonicalización de aliases.
\end{itemize}

\begin{center}\rule{0.5\linewidth}{0.5pt}\end{center}

\hypertarget{arquitectura-recomendada-para-producciuxf3n}{%
\subsection{8. Arquitectura recomendada para
producción}\label{arquitectura-recomendada-para-producciuxf3n}}

\begin{Shaded}
\begin{Highlighting}[]
\NormalTok{┌──────────────────────────────┐}
\NormalTok{│ Wallet o aplicación cliente │}
\NormalTok{└──────────────┬───────────────┘}
\NormalTok{               │}
\NormalTok{               ▼}
\NormalTok{┌──────────────────────────────┐}
\NormalTok{│ SatsPath SDK                 │}
\NormalTok{│                              │}
\NormalTok{│ {-} Resolver chain             │}
\NormalTok{│ {-} Verificación               │}
\NormalTok{│ {-} Trust model                │}
\NormalTok{│ {-} Router y scoring           │}
\NormalTok{│ {-} Payload builder            │}
\NormalTok{└───────┬───────────┬──────────┘}
\NormalTok{        │           │}
\NormalTok{        │           └─────────────────────┐}
\NormalTok{        ▼                                 ▼}
\NormalTok{┌───────────────┐               ┌──────────────────┐}
\NormalTok{│ Resolvers     │               │ Wallet adapter   │}
\NormalTok{│               │               │                  │}
\NormalTok{│ DNS / HTTPS   │               │ Lightning        │}
\NormalTok{│ Nostr / P2P   │               │ Bitcoin on{-}chain │}
\NormalTok{│ Platform      │               │ Ark              │}
\NormalTok{└───────────────┘               └──────────────────┘}
\end{Highlighting}
\end{Shaded}

La integración ideal mantiene una separación estricta:

\hypertarget{satspath-core}{%
\subsubsection{SatsPath Core}\label{satspath-core}}

Responsable de:

\begin{itemize}
\tightlist
\item
  Tipos.
\item
  Firmas.
\item
  Verificación.
\item
  Resolución.
\item
  Routing.
\item
  Construcción de instrucciones.
\end{itemize}

\hypertarget{wallet-adapter}{%
\subsubsection{Wallet Adapter}\label{wallet-adapter}}

Responsable de:

\begin{itemize}
\tightlist
\item
  Obtener balance.
\item
  Consultar capacidades.
\item
  Confirmar con el usuario.
\item
  Pagar invoice.
\item
  Crear y firmar transacción.
\item
  Reportar éxito o fallo.
\end{itemize}

\hypertarget{transport-adapter}{%
\subsubsection{Transport Adapter}\label{transport-adapter}}

Responsable de:

\begin{itemize}
\tightlist
\item
  DNS.
\item
  HTTPS.
\item
  Nostr.
\item
  P2P.
\item
  Plataforma.
\end{itemize}

\begin{center}\rule{0.5\linewidth}{0.5pt}\end{center}

\hypertarget{primera-versiuxf3n-de-producciuxf3n-recomendada}{%
\subsection{9. Primera versión de producción
recomendada}\label{primera-versiuxf3n-de-producciuxf3n-recomendada}}

La primera versión no necesita soportar todas las funciones del
repositorio.

Debe enfocarse en:

\begin{Shaded}
\begin{Highlighting}[]
\NormalTok{Resolve}
\NormalTok{Verify}
\NormalTok{Route}
\NormalTok{Generate payable instruction}
\NormalTok{Open external wallet}
\end{Highlighting}
\end{Shaded}

\hypertarget{alcance-recomendado}{%
\subsubsection{Alcance recomendado}\label{alcance-recomendado}}

\begin{itemize}
\tightlist
\item
  Lightning Address y LNURL-pay.
\item
  Invoice BOLT11 validada.
\item
  Dirección on-chain y BIP-21.
\item
  BIP-353 para dominios controlados.
\item
  HTTPS \texttt{.well-known}.
\item
  Perfil firmado y expiración.
\item
  QR.
\item
  Wallet handoff.
\item
  Ark como adapter experimental interoperable.
\end{itemize}

\hypertarget{fuera-del-primer-alcance}{%
\subsubsection{Fuera del primer
alcance}\label{fuera-del-primer-alcance}}

\begin{itemize}
\tightlist
\item
  Split payments complejos.
\item
  Ejecución automática multirail.
\item
  Recuperación social.
\item
  Registro global propio.
\item
  Todos los proveedores DNS.
\item
  Todos los runtimes P2P.
\end{itemize}

Reducir el alcance no disminuye la visión. Evita que una arquitectura
prometedora se convierta en una colección de módulos a medio terminar,
que es el destino natural de demasiados proyectos técnicamente
ambiciosos.

\begin{center}\rule{0.5\linewidth}{0.5pt}\end{center}

\hypertarget{criterios-de-aceptaciuxf3n}{%
\subsection{10. Criterios de
aceptación}\label{criterios-de-aceptaciuxf3n}}

SatsPath debería considerarse listo para una integración productiva
inicial cuando cumpla:

\hypertarget{seguridad}{%
\subsubsection{Seguridad}\label{seguridad}}

\begin{itemize}
\tightlist
\item[$\square$]
  Los perfiles inválidos son rechazados.
\item[$\square$]
  Los perfiles expirados son rechazados.
\item[$\square$]
  Los replays con menor secuencia son rechazados.
\item[$\square$]
  Las claves de identidad están protegidas por el sistema operativo o
  cifrado.
\item[$\square$]
  La API administrativa requiere autenticación.
\item[$\square$]
  No existe CORS abierto en una instancia pública.
\item[$\square$]
  Los fetches remotos tienen protección SSRF.
\end{itemize}

\hypertarget{identidad}{%
\subsubsection{Identidad}\label{identidad}}

\begin{itemize}
\tightlist
\item[$\square$]
  La procedencia del perfil se incluye en la respuesta.
\item[$\square$]
  El cliente distingue firma de perfil e identidad verificada.
\item[$\square$]
  BIP-353 funciona con DNSSEC.
\item[$\square$]
  Los challenges de plataforma son reales.
\end{itemize}

\hypertarget{pagos}{%
\subsubsection{Pagos}\label{pagos}}

\begin{itemize}
\tightlist
\item[$\square$]
  Una invoice LNURL válida puede abrirse y pagarse desde una wallet
  externa.
\item[$\square$]
  Una invoice con monto incorrecto es rechazada.
\item[$\square$]
  Una invoice expirada es rechazada.
\item[$\square$]
  El URI BIP-21 es compatible con al menos dos wallets.
\item[$\square$]
  El adapter Ark ha sido probado con una implementación concreta.
\end{itemize}

\hypertarget{integraciuxf3n}{%
\subsubsection{Integración}\label{integraciuxf3n}}

\begin{itemize}
\tightlist
\item[$\square$]
  Existe una API estable y versionada.
\item[$\square$]
  Existe una librería usable en al menos una plataforma cliente.
\item[$\square$]
  Hay test vectors públicos.
\item[$\square$]
  Existe una aplicación demo que utiliza el SDK, no lógica duplicada.
\end{itemize}

\hypertarget{operaciones}{%
\subsubsection{Operaciones}\label{operaciones}}

\begin{itemize}
\tightlist
\item[$\square$]
  CI es obligatorio.
\item[$\square$]
  Releases están firmadas.
\item[$\square$]
  Existe changelog.
\item[$\square$]
  Hay métricas sin información sensible.
\item[$\square$]
  Existe política de vulnerabilidades.
\end{itemize}

\begin{center}\rule{0.5\linewidth}{0.5pt}\end{center}

\hypertarget{roadmap-sugerido}{%
\subsection{11. Roadmap sugerido}\label{roadmap-sugerido}}

\hypertarget{fase-1-consolidaciuxf3n-del-core}{%
\subsection{Fase 1: consolidación del
core}\label{fase-1-consolidaciuxf3n-del-core}}

\begin{itemize}
\tightlist
\item
  Limpiar documentación contradictoria.
\item
  Cerrar o actualizar issues obsoletos.
\item
  Convertir CI en requisito obligatorio.
\item
  Definir contrato \texttt{ResolvedProfile}.
\item
  Mejorar estados de error.
\item
  Aplicar políticas de replay y secuencia.
\end{itemize}

\hypertarget{fase-2-satspath-preview}{%
\subsection{Fase 2: SatsPath Preview}\label{fase-2-satspath-preview}}

\begin{itemize}
\tightlist
\item
  Perfil real publicado en HTTPS.
\item
  BIP-353 funcional.
\item
  LNURL y BOLT11 interoperables.
\item
  BIP-21 interoperable.
\item
  QR y wallet handoff.
\item
  PWA basada en el SDK.
\end{itemize}

\hypertarget{fase-3-seguridad-operativa}{%
\subsection{Fase 3: seguridad
operativa}\label{fase-3-seguridad-operativa}}

\begin{itemize}
\tightlist
\item
  Protección de claves.
\item
  Auth del daemon.
\item
  CORS seguro.
\item
  SSRF protection.
\item
  Rate limiting.
\item
  Auditoría de dependencias.
\end{itemize}

\hypertarget{fase-4-integraciuxf3n-de-wallet}{%
\subsection{Fase 4: integración de
wallet}\label{fase-4-integraciuxf3n-de-wallet}}

\begin{itemize}
\tightlist
\item
  Elegir una wallet de referencia.
\item
  Integrar WASM, FFI o Rust directamente.
\item
  Aportar capacidades de la wallet al router.
\item
  Reportar resultado del pago.
\item
  Implementar fallback después de un fallo.
\end{itemize}

\hypertarget{fase-5-ark-y-rails-adicionales}{%
\subsection{Fase 5: Ark y rails
adicionales}\label{fase-5-ark-y-rails-adicionales}}

\begin{itemize}
\tightlist
\item
  Ark interoperable.
\item
  BOLT12 completo.
\item
  Silent Payments.
\item
  Split Payments con semántica definida.
\end{itemize}

\hypertarget{fase-6-especificaciuxf3n-abierta}{%
\subsection{Fase 6: especificación
abierta}\label{fase-6-especificaciuxf3n-abierta}}

\begin{itemize}
\tightlist
\item
  Congelar SatsPath Protocol v1.
\item
  Publicar test vectors.
\item
  Crear suite de conformidad.
\item
  Implementar un segundo cliente independiente.
\item
  Solicitar revisión externa.
\end{itemize}

\begin{center}\rule{0.5\linewidth}{0.5pt}\end{center}

\hypertarget{posicionamiento-recomendado}{%
\subsection{12. Posicionamiento
recomendado}\label{posicionamiento-recomendado}}

SatsPath debería presentarse como:

\begin{quote}
Un protocolo abierto y no custodial para resolver identidades de pago
Bitcoin, verificar perfiles firmados y seleccionar la mejor instrucción
disponible entre múltiples rails.
\end{quote}

No debería limitarse a describirse como:

\begin{itemize}
\tightlist
\item
  Una wallet.
\item
  Un explorador de direcciones.
\item
  Un registro de usuarios.
\item
  Un router basado solamente en fees.
\item
  Un sistema P2P específico.
\end{itemize}

Su valor está en unir cuatro capas:

\begin{Shaded}
\begin{Highlighting}[]
\NormalTok{Identidad}
\NormalTok{Verificación}
\NormalTok{Descubrimiento}
\NormalTok{Routing}
\end{Highlighting}
\end{Shaded}

\begin{center}\rule{0.5\linewidth}{0.5pt}\end{center}

\hypertarget{conclusiuxf3n}{%
\subsection{13. Conclusión}\label{conclusiuxf3n}}

SatsPath ya posee una base funcional considerable:

\begin{itemize}
\tightlist
\item
  Perfiles firmados.
\item
  Validación criptográfica.
\item
  Resolución multitransporte.
\item
  LNURL y BOLT11.
\item
  BIP-353.
\item
  Nostr.
\item
  P2P.
\item
  Router.
\item
  QR y wallet handoff.
\item
  WASM y bindings móviles.
\end{itemize}

También puede facilitar pagos reales porque trabaja con información
pública válida y produce instrucciones pagables por una wallet. No
necesita custodiar fondos para tener utilidad práctica.

El trabajo restante no consiste principalmente en agregar más
protocolos. Consiste en endurecer la seguridad, demostrar
interoperabilidad y convertir las piezas actuales en una experiencia
estable que otra wallet pueda integrar sin depender de supuestos
internos del repositorio.

El punto que marcará la transición de proyecto a infraestructura será el
siguiente:

\begin{quote}
Dos wallets independientes resuelven el mismo identificador, verifican
el mismo perfil, seleccionan una ruta compatible y completan el pago
utilizando la instrucción producida por SatsPath.
\end{quote}

Cuando ese flujo sea reproducible, versionado y seguro, SatsPath dejará
de ser únicamente una implementación prometedora y comenzará a funcionar
como un protocolo de interoperabilidad real para pagos Bitcoin.

\begin{center}\rule{0.5\linewidth}{0.5pt}\end{center}

\hypertarget{referencias-dentro-del-repositorio}{%
\subsection{14. Referencias dentro del
repositorio}\label{referencias-dentro-del-repositorio}}

\begin{Shaded}
\begin{Highlighting}[]
\NormalTok{README.md}
\NormalTok{docs/protocol.md}
\NormalTok{docs/resolvers.md}
\NormalTok{docs/implementations.md}
\NormalTok{docs/threat\_model.md}
\NormalTok{docs/wire\_p2p.md}
\NormalTok{crates/satspath{-}core/src/}
\NormalTok{crates/satspath{-}router/src/}
\NormalTok{crates/satspath{-}cli/src/}
\NormalTok{crates/satspathd/src/main.rs}
\NormalTok{crates/satspath{-}wasm/src/}
\NormalTok{crates/satspath{-}ffi/src/}
\NormalTok{sdk/satspath{-}p2p/}
\NormalTok{.github/workflows/ci.yml}
\end{Highlighting}
\end{Shaded}


\end{document}
